% =============================================================
% 05_foundation.tex — 基礎理論・ハードウェア・ソフトウェア
% =============================================================
\chapter{基礎理論・ハードウェア・ソフトウェア}

\chaptermeta{%
  \item 10進数と2進数の相互変換ができる
  \item AND/OR/NOT/XORの論理演算を使える
  \item CPU・メモリの階層構造を説明できる
  \item RAID 0/1/5の違いを即答できる
  \item OSの役割とOSSの特徴を整理できる
}{90分}{3}{第1--4章}

テクノロジ系の最初は「コンピュータの基礎」です。
2進数・論理演算は\textbf{計算問題}として出題され、
CPU/メモリ/RAIDは\textbf{用語の識別問題}として頻出します。

% =============================================================
\section{数の表現}
% =============================================================

\begin{teigi}[基数変換の基本]
\begin{itemize}
  \item \termtag{10進数→2進数}：2で割り続け、余りを下から読む
  \item \termtag{2進数→10進数}：各桁に位の重み（$2^0, 2^1, 2^2, \ldots$）を掛けて足す
  \item \termtag{2進数→16進数}：4桁ずつ区切って16進に変換
  \item 16進数は先頭に「0x」や末尾に「H」をつけて表記（例：0x1F、1FH）
\end{itemize}
\end{teigi}

% --- 例題 ---
\begin{reidai}{10進数から2進数への変換}
10進数の13を2進数で表したものはどれか。

\sentakushi{1011}{1100}{1101}{1110}
\end{reidai}

\begin{shishin}
2で割っていく。$13 \div 2 = 6 \cdots 1$、$6 \div 2 = 3 \cdots 0$、$3 \div 2 = 1 \cdots 1$、$1 \div 2 = 0 \cdots 1$。
余りを下から読むと$1101$。
\end{shishin}

\begin{kaitou}
\textbf{正解：ウ（1101）}

$13 = 8 + 4 + 1 = 2^3 + 2^2 + 2^0 = 1101_{(2)}$。
別解として位の重みで分解する方法でも確認できる。
\end{kaitou}

\begin{minuki}[見抜き方]
\begin{itemize}
  \item 小さい数（0〜15）は暗記してしまうと速い：$1010=10, 1100=12, 1101=13, 1111=15$
  \item 2の累乗を覚える：$1, 2, 4, 8, 16, 32, 64, 128, 256, 512, 1024$
  \item 検算は2進数→10進数変換で行う
\end{itemize}
\end{minuki}

\subsection{補数と負の数}

\begin{teigi}[2の補数]
コンピュータ内部で負の数を表す方式。
\begin{enumerate}
  \item 全ビットを反転（0→1、1→0）
  \item 結果に1を足す
\end{enumerate}
例：$0101_{(2)}$（＝5）の2の補数 → $1010 + 1 = 1011_{(2)}$（＝$-5$）。
\end{teigi}

% =============================================================
\section{論理演算}
% =============================================================

\begin{hikaku}[4つの基本論理演算]
\comptable{演算 & 記号 & 結果が1になる条件}{
AND（論理積） & $\cdot$ & 両方1のときだけ \\
OR（論理和） & $+$ & どちらか一方でも1 \\
NOT（否定） & $\overline{A}$ & 入力の反転 \\
XOR（排他的論理和） & $\oplus$ & 2つが異なるとき
}
\end{hikaku}

\begin{pitfall}[ORとXORの混同]
\begin{itemize}
  \item OR：「1 OR 1 = 1」（どちらかが1ならOK）
  \item XOR：「1 XOR 1 = 0」（同じなら0、違えば1）
  \item 試験では「排他的」というキーワードでXORを識別
\end{itemize}
\end{pitfall}

% =============================================================
\section{アルゴリズムの基礎}
% =============================================================

\begin{teigi}[基本的なアルゴリズム]
\begin{itemize}
  \item \termtag{線形探索}：先頭から順に探す（単純だがデータ量に比例して遅い）
  \item \termtag{2分探索}：ソート済みデータを半分ずつ絞る（高速：$O(\log n)$）
  \item \termtag{バブルソート}：隣り合うデータを比較・交換して並べ替え
  \item \termtag{選択ソート}：最小値を見つけて先頭と交換を繰り返す
\end{itemize}
\end{teigi}

\begin{pitfall}[「2分探索はソート不要」は誤り]
2分探索の前提は\textbf{データがソート済み}であること。
未整列データに2分探索を適用しても正しい結果は得られない。
\end{pitfall}

% =============================================================
\section{CPU・メモリの階層構造}
% =============================================================

\begin{teigi}[記憶装置の階層]
速度が速い順（＝容量が小さい順）：
\begin{enumerate}
  \item \termtag{レジスタ}：CPU内部の超高速記憶（数十〜数百個）
  \item \termtag{キャッシュメモリ}：CPUとメインメモリの間の高速バッファ（SRAM）
  \item \termtag{メインメモリ（主記憶）}：実行中プログラムが置かれる（DRAM）
  \item \termtag{SSD/HDD（補助記憶）}：電源を切ってもデータが残る
\end{enumerate}
上に行くほど\textbf{高速・小容量・高価格}、下に行くほど\textbf{低速・大容量・低価格}。
\end{teigi}

\begin{hikaku}[SRAMとDRAM]
\comptable{特性 & SRAM & DRAM}{
速度 & 高速 & 低速 \\
容量あたりの価格 & 高い & 安い \\
リフレッシュ & 不要 & 必要 \\
用途 & キャッシュメモリ & メインメモリ
}
\end{hikaku}

\subsection{CPUの性能指標}

\begin{teigi}[クロック周波数とCPI]
\begin{itemize}
  \item \termtag{クロック周波数}：CPUが1秒間に刻むクロック数（単位：Hz）。大きいほど高速。
  \item \termtag{MIPS}（Million Instructions Per Second）：1秒間に実行できる命令数（百万単位）
  \item クロック周波数1GHz、1命令に2クロック必要 → MIPS $= 1{,}000 \div 2 = 500$ MIPS
\end{itemize}
\end{teigi}

% =============================================================
\section{RAID}
% =============================================================

\begin{hikaku}[RAID 0/1/5の比較]
\comptable{RAID & 方式 & 特徴}{
RAID 0 & ストライピング & 高速だが冗長性なし（1台故障で全滅） \\
RAID 1 & ミラーリング & 同じデータを2台に書く（容量半減） \\
RAID 5 & パリティ分散 & 3台以上、1台故障まで復旧可能
}
\end{hikaku}

\begin{pitfall}[RAID 5の「1台まで」を忘れる]
RAID 5は\textbf{1台}の故障まで復旧可能。2台同時故障ではデータが失われる。
2台故障に耐えるのはRAID 6（パリティを二重化）。
\end{pitfall}

% =============================================================
\section{OS・OSS}
% =============================================================

\begin{teigi}[OSの役割]
\begin{itemize}
  \item \termtag{OS}（Operating System）：ハードウェアを抽象化し、アプリケーションに共通の実行環境を提供
  \item 主な機能：プロセス管理、メモリ管理、ファイル管理、入出力管理
  \item \termtag{仮想記憶}：主記憶が不足したとき、補助記憶の一部を主記憶のように使う仕組み
\end{itemize}
\end{teigi}

\begin{teigi}[OSSの特徴]
\termtag{OSS}（Open Source Software）：ソースコードが公開され、自由に利用・改変・再配布できるソフトウェア。
\begin{itemize}
  \item 無償とは限らない（サポートは有償の場合あり）
  \item 代表例：Linux、Apache、MySQL、PostgreSQL、Python、Firefox
  \item ライセンス例：GPL（改変後も公開義務）、MIT（制約が少ない）
\end{itemize}
\end{teigi}

\begin{pitfall}[「OSSは無料」は不正確]
OSSは\textbf{ソースコードが公開}されていることが本質。
商用サポートを有償で提供することは可能であり、OSS＝無料ではない。
\end{pitfall}

% =============================================================
\section{過去問ドリル}
% =============================================================

\shoumatsuhead{過去問ドリル：基礎理論とハードウェア}

\begin{itpmon}[\sourcebadge{original}\enspace\diffbadge{標}]{%
RAID 5の説明として、最も適切なものはどれか。}
\sentakushi
  {データを複数のディスクに分散して書き込み、高速化するが冗長性はない。}
  {同一データを2台のディスクに書き込み、1台が故障してもデータを保持できる。}
  {3台以上のディスクにデータとパリティを分散し、1台の故障まで復旧可能である。}
  {2台のディスクにそれぞれ独立したパリティを持ち、2台の故障まで復旧可能である。}
\end{itpmon}
\begin{kaisetsu}{ウ}
\textbf{RAID 5}はデータとパリティを3台以上に分散記録する。アはRAID 0、イはRAID 1、エはRAID 6の説明。
\end{kaisetsu}

\begin{itpmon}[\sourcebadge{original}\enspace\diffbadge{標}]{%
CPUのクロック周波数に関する説明として、最も適切なものはどれか。}
\sentakushi
  {クロック周波数が高いほど、1秒あたりに実行できる命令数が多くなる。}
  {クロック周波数が高いほど、記憶装置の容量が大きくなる。}
  {クロック周波数は、ネットワークの通信速度を表す指標である。}
  {クロック周波数は、ディスプレイの解像度を決定する要素である。}
\end{itpmon}
\begin{kaisetsu}{ア}
\textbf{クロック周波数}はCPUの動作速度の指標。周波数が高いほど単位時間に多くの命令を処理できる。記憶容量・通信速度・解像度とは直接関係しない。
\end{kaisetsu}

\begin{itpmon}[\sourcebadge{original}\enspace\diffbadge{標}]{%
OSSに関する説明として、最も適切なものはどれか。}
\sentakushi
  {ソースコードが非公開のソフトウェアである。}
  {利用にはメーカーのライセンス契約が必須である。}
  {ソースコードが公開され、自由に利用・改変・再配布が認められている。}
  {企業が開発したソフトウェアの体験版である。}
\end{itpmon}
\begin{kaisetsu}{ウ}
\textbf{OSS}はソースコードの公開が本質。GPL、MIT、Apacheライセンスなど、様々なライセンス形態がある。自由に利用・改変・再配布が認められている。
\end{kaisetsu}

% =============================================================
\section{タイムアタック}
% =============================================================

\begin{timeattack}{基礎理論クイックチェック}{6分}{5問中4問}

\begin{itpmon}[\diffbadge{基}]{10進数の10を2進数で表したものはどれか。}
\sentakushi{1000}{1001}{1010}{1011}
\end{itpmon}
\begin{kaisetsu}{ウ}
$10 = 8 + 2 = 2^3 + 2^1 = 1010_{(2)}$。
\end{kaisetsu}

\begin{itpmon}[\diffbadge{基}]{論理演算「1 AND 0」の結果はどれか。}
\sentakushi{0}{1}{2}{不定}
\end{itpmon}
\begin{kaisetsu}{ア}
AND（論理積）は両方1のときだけ1。$1 \text{ AND } 0 = 0$。
\end{kaisetsu}

\begin{itpmon}[\diffbadge{標}]{キャッシュメモリの説明として正しいものはどれか。}
\sentakushi
  {電源を切ってもデータが残る記憶装置}
  {CPUとメインメモリの間に置かれる高速な記憶装置}
  {大容量で低価格な磁気ディスク装置}
  {プログラムの実行中にデータを一時的に退避する補助記憶装置}
\end{itpmon}
\begin{kaisetsu}{イ}
\textbf{キャッシュメモリ}はCPUとメインメモリの速度差を埋めるSRAM。アはSSD/HDD、ウもHDD、エはスワップ領域の説明。
\end{kaisetsu}

\begin{itpmon}[\diffbadge{基}]{RAID 1の方式名はどれか。}
\sentakushi{ストライピング}{ミラーリング}{パリティ分散}{二重パリティ}
\end{itpmon}
\begin{kaisetsu}{イ}
RAID 1 = \textbf{ミラーリング}。アはRAID 0、ウはRAID 5、エはRAID 6。
\end{kaisetsu}

\begin{itpmon}[\diffbadge{標}]{仮想記憶の説明として正しいものはどれか。}
\sentakushi
  {CPUの内部にある超高速な記憶領域}
  {主記憶の不足分を補助記憶で補い、見かけ上の主記憶を拡大する仕組み}
  {データを複数のディスクに分散する技術}
  {ネットワーク上の記憶領域を共有する仕組み}
\end{itpmon}
\begin{kaisetsu}{イ}
\textbf{仮想記憶}は主記憶が不足したとき、補助記憶（SSD/HDD）の一部を使って見かけ上の主記憶容量を拡大する技術。
\end{kaisetsu}

\end{timeattack}

% =============================================================
\section{よくあるミスと到達確認}
% =============================================================

\begin{yokumiss}[この章でよくあるミス]
\begin{enumerate}[leftmargin=1.6em, itemsep=5pt]
  \item \textbf{2進数変換で余りの読む方向を間違える}——下から上に読む。
  \item \textbf{ORとXORを混同する}——OR: 1 OR 1 = 1、XOR: 1 XOR 1 = 0。
  \item \textbf{SRAMとDRAMの用途を逆にする}——SRAM=キャッシュ、DRAM=メインメモリ。
  \item \textbf{RAID 0に冗長性があると思い込む}——RAID 0は高速化のみ、冗長性なし。
  \item \textbf{OSSを「無料」と同義に扱う}——本質はソースコード公開。商用サポートは有償可。
\end{enumerate}
\end{yokumiss}

\begin{tatsaku}
  \item 10進数と2進数を相互変換できる
  \item AND/OR/NOT/XORの真理値表を書ける
  \item 記憶装置の階層（レジスタ→キャッシュ→メインメモリ→補助記憶）を言える
  \item SRAMとDRAMの違いを3点以上挙げられる
  \item RAID 0/1/5の違いを即答できる
  \item 仮想記憶の仕組みを説明できる
  \item OSSの定義と代表例を3つ以上言える
  \item GPLとMITライセンスの違いを説明できる
\end{tatsaku}

\nextchapter{%
ネットワークとデータベースに進みます。TCP/IP、プロトコル、SQL、正規化——
「つながる仕組み」と「データの管理」を学びます。
}
